\section{Light vectors and the retained spectrum}\label{sec:vectors}

\subsection{The physical vector operator}
On real skew-adjoint generators set
$\BF(U,V)=-\tfrac12\Tr_{32}(UV)$. Define $g_{10}$ by
\begin{equation}\label{eq:YMaction}
 S_{\rm YM}=-\frac1{4g_{10}^2}\int
    \BF(F_{MN},F^{MN})\,d\Vol_{10}.
\end{equation}
The constant dilaton is absorbed in this coupling. Initially use the
product spacetime metric $g_{1,5}+L^2g_0$; a constant change to
six-dimensional Einstein frame must rescale masses and scalar kinetic
tensors consistently. On a realified unitary block $\BF$ restricts to
$-\operatorname{Re}\Tr$. A multiplicity generator at $E_0$ consequently
has reference norm $3V_0\tr_2(\xi^*\xi)$.

For a Lorentz-transverse external fluctuation
$\delta\mathcal A_\mu=v_\mu\psi(y)$, background gauge separates the
internal Goldstone variation. The external kinetic and mass forms are
$g_{10}^{-2}\|\psi\|^2_{L^2(L^2g_0)}$ and
$g_{10}^{-2}\|d_{\mathcal A}\psi\|^2_{L^2(L^2g_0)}$.
The coupling cancels in their generalized eigenvalue. On endomorphism
zero-forms the HE K\"ahler identity gives
\begin{equation}\label{eq:physicaloperator}
 d_{\mathcal A}^*d_{\mathcal A}=2\Delta_{\dbar,\End E}.
\end{equation}
The induced endomorphism mean curvature is zero, so the $\partial$
and $\dbar$ Laplacians agree in this degree. This is a
zero-form vector-profile identity, not an identity for the full
one-form Yang--Mills Hessian; compare \cite[appendix~A]{louis}.
The zero-form and one-form norms scale as $L^4$ and $L^2$,
respectively, so physical squared masses acquire the factor $L^{-2}$.

For $\xi=\sum z_j\sigma_j$ the finite Dolbeault energy and mass are
\begin{align}
 q_x(\xi)&=\sum_j\|[\xi,X_j]\|^2+\|[\xi,Y_j]\|^2
 =8\bigl(a^2|z_1|^2+(a^2+b^2)|z_2|^2+b^2|z_3|^2\bigr),
 \label{eq:pauli-energy}\\
 M_0(\xi)&=3V_0\tr(\xi^*\xi)=6V_0\sum_j|z_j|^2.
 \label{eq:pauli-mass}
\end{align}
Thus its ordered eigenvalues for $0<b\le a$ are
$\tfrac4{3V_0}(b^2,a^2,a^2+b^2)$.

\begin{proposition}\label{prop:vectors}
For the actual bundles of \eqref{eq:pauli}, sufficiently near the
origin, there are exactly three real vector species whose positive
masses tend to zero away from the unbroken kernel. They satisfy
\eqref{eq:intro-masses}, with a constant independent of $b/a$.
The corresponding two-sided inequality retains exact zeros on the
axes and at the origin. All other nonzero vector-profile eigenvalues
have a fixed positive lower bound in reference units.
\end{proposition}
\begin{proof}
Apply \cite[Theorem~1.1(ii)]{wuebben15} at the locally free polystable
bundle $A\otimes\CC^2$, at constant polarization and fixed
determinant. Its trace-free multiplicity algebra has dimension three.
The theorem bounds every small Dolbeault eigenvalue between
$(1-C\|x\|^2)\nu_j(x)$ and $(1+C\|x\|^2)\nu_j(x)$,
including zeros, and gives a positive complementary gap. Here
$\|x\|^2=2(a^2+b^2)$ and \eqref{eq:pauli-energy}--\eqref{eq:pauli-mass}
give $\nu_j$. Equation~\eqref{eq:physicaloperator} and the internal
scale give \eqref{eq:intro-masses}. The real adjoint sector is preserved
by this Laplacian. Each complex eigenspace is the complexification of
its Hermitian, or equivalently skew-Hermitian, real form. Its complex
multiplicity is therefore the number of real generator profiles,
not twice that number. The remaining orthogonal-adjoint blocks have
the kernels already identified in \eqref{eq:centralizer}; their
positive spectrum stays separated from zero by elliptic continuity
on the fixed compact background.
\end{proof}

The essential input is the relative form of the imported estimate.
In its proof the complementary Schur term is bounded by the finite
energy of the same vector. A uniform additive error of order
$(a^2+b^2)^2$ would not establish the result when $b\ll a^2$.
There is no additional gauge coupling multiplying the masses in
\eqref{eq:intro-masses} after the convention \eqref{eq:YMaction} is fixed.

\subsection{Massless fields and Higgs multiplets}
The vector gap does not make the remaining matter heavy. The full
orthogonal adjoint decomposes into endomorphism, exterior-square,
mixed tensor and fourteen-dimensional vector sectors. For a stable
determinant-trivial rank-$r$ bundle $F$ with $c_2(F)=k$ and the
Hom vanishings established above, Riemann--Roch and Serre duality give
\begin{align}\label{eq:cohomology}
 h^1(F)&=k-2r,&
 \dim_{\mathbb H}\mathcal M_F&=rk-r^2+1,\\
 h^1(F\otimes G)&=r_Gk_F+r_Fk_G-2r_Fr_G,&
 h^1(\Lambda^2F)&=(r-2)k-r(r-1).\nonumber
\end{align}
These follow by integrating $\operatorname{ch}(\cdot)\operatorname{td}(S)$;
the exterior-square formula uses
$c_2(\Lambda^2F)=(r-2)c_2(F)$ when $c_1(F)=0$.
The neutral moduli formula follows from $h^1(\End F)=2rk-2r^2+2$.
In the spectrum convention of \cite{honecker}, complex conjugate
sectors combine into full hypermultiplets. At a generic point the
dimension-weighted counts are listed in table~\ref{tab:matter}.

\begin{table}[t]
\centering
\caption{Generic massless hypermultiplet contributions. Charges refer to
the two central $\U(1)$ factors. The fourteen-dimensional sectors contain
hypermultiplets in the vector representation of $\SO(14)$, not additional
gauge vectors.}\label{tab:matter}
\begin{adjustbox}{max width=\textwidth,center}
\begin{tabular}{lr}
\toprule
Sector & Dimension-weighted hypermultiplet count\\
\midrule
$E$ moduli, neutral &49\\
$B$ moduli, neutral &22\\
$\Lambda^2E$, charge $(2,0)$ &26\\
$\Lambda^2B$, charge $(0,2)$ &4\\
$E\otimes B$, charge $(1,1)$ &66\\
$E\otimes B^*$, charge $(1,-1)$ &66\\
$E$ sector in $\mathbf {14}$, charge $(1,0)$ &28\\
$B$ sector in $\mathbf {14}$, charge $(0,1)$ &56\\
K3 geometry and two-form fields &20\\
\midrule
Total &337\\
\bottomrule
\end{tabular}
\end{adjustbox}
\end{table}

At enhancement the $E$ endomorphism sector consists of thirteen
adjoint-$\U(2)$ hypermultiplets and contributes 52. On a nonzero axis
the two distinct rank-three factors contribute $2\cdot13$ neutral
hypermultiplets and 24 off-diagonal charged ones, giving 50.
The other dimension-weighted totals are unchanged. Thus
\begin{equation}\label{eq:counts}
 (n_H,n_V)=
 \begin{cases}
 (340,96),&a=b=0,\\
 (338,94),&ab=0,\quad a^2+b^2>0,\\
 (337,93),&ab\ne0.
 \end{cases}
\end{equation}
There is one tensor multiplet, containing the six-dimensional dilaton,
and the gravity multiplet. Every row obeys
$n_H-n_V+29n_T=273$. This is a gravitational anomaly count; it does
not replace the gauge-anomaly and embedding analysis of \cite{honecker}.
As a separate index check, the real gauge ASD complex has
$h^1-h^0-h^2=4\cdot30\cdot24-4\cdot496=896$ on K3. Our generic
gauge hypermultiplets contribute $4(337-20)$, and $h^0=93$ gives
$4(337-20)-4\cdot93=896$; the corresponding enhanced and axis
counts agree as well. Here the expression subtracts all four
parallel contributions, equivalently using $h^2=3h^0$ in the
ASD index.

Three hypermultiplets and three vectors combine into three massive
$\mathcal N=(1,0)$ vector multiplets generically. Each massive vector
has five physical polarizations and three physical real scalar
partners; the fourth scalar of its hypermultiplet is a Goldstone
mode. On an axis only two such multiplets are massive.

\subsection{The common low-energy space}
Choose background gauge, diffeomorphism gauge and two-form gauge at
the reference background. The leading internal operators are the
adjoint zero-form Laplacian for vectors, the positive ASD deformation
square for connection scalars, the physical Lichnerowicz operator for
metric deformations, ordinary Hodge Laplacians for forms, and the
corresponding Dirac complexes for fermions. The ASD square is
$d_{\mathcal A}d_{\mathcal A}^*+2d_{\mathcal A}^{+*}d_{\mathcal A}^+$
on one-forms. The parallel-spinor factorization gives nonnegativity
of the physical Ricci-flat metric operator \cite{parallelspinor};
this excludes the unphysical Euclidean conformal mode. Compactness
gives finite kernels and positive complementary gaps. Supersymmetry
identifies the physical superpartners.

Take a cutoff below the minimum of these reference gaps and retain
every reference low mode, including the gauge/Goldstone complex.
Smooth elliptic dependence gives common finite low spectral clusters
and a complementary gap $\gamma>0$ for all sufficiently small
$(a,b)$. One must not first divide out a different number of scalar
gauge directions at each stratum and then assert continuity on a
fixed space. The theory with the Higgs multiplets retained has the
common leading window $E_{\rm probe}\ll\sqrt\gamma/L$.
Eliminating them instead requires $E_{\rm probe}\ll b/L$ on a
generic regular patch. Higher-derivative string terms are used
perturbatively through leading resolvents, not as a positive
second-order operator on the entire tower.
